\begin{tabularx}{\textwidth}{@{}>{\raggedright\arraybackslash}p{0.20\textwidth} X@{}}
\toprule
Feature & Description \\
\midrule
Topical responsiveness & How directly a reply addresses the poster's actual view rather than pivoting to meta-issues or adjacent topics. \\
Premise targeting & Whether a reply attacks the core assumption or inference supporting the view instead of focusing on side details. \\
Clarification demand & How much a reply advances the discussion by asking the poster to specify vague claims, terms, or examples. \\
Conceptual reframing & Whether a reply changes the frame of the debate by redefining the key concept or explaining that the dispute is really about something else. \\
Thought experiment & Whether a reply uses an invented scenario to test the implications or consistency of the poster's view. \\
Analogy use & How much a reply explains its point by mapping the issue onto a different familiar domain. \\
Consistency test & Whether a reply probes the poster's logic by applying it to parallel cases and checking whether the result is acceptable. \\
Reductio extension & Whether a reply undermines a view by extending it to an absurd or unacceptable conclusion. \\
Counterexample case & Whether a reply primarily challenges a general claim by presenting a concrete case that does not fit it. \\
Firsthand testimony & How much a reply relies on the writer's own experiences or preferences as evidence or illustration. \\
Stakeholder perspective & Whether a reply argues from direct membership in the affected group or a role that gives special practical exposure to the issue. \\
Empirical citation & How much a reply supports its claims with studies, data, or directly linked factual evidence. \\
Expert consensus appeal & Whether a reply leans on the agreement of qualified experts as a reason to update the view. \\
Legal mechanics & How much a reply explains the governing law, procedure, or institutional rules behind the issue. \\
Implementation realism & Whether a reply evaluates how workable a proposal would be once logistics, administration, and real-world constraints are considered. \\
Economic reasoning & How much a reply argues in terms of prices, incentives, resources, employment, or market effects. \\
Quantified argument & Whether a reply strengthens its case by attaching numerical magnitude or rates to the claims it makes. \\
Causal mechanism & Whether a reply explains the process by which one thing is supposed to produce another, rather than merely asserting a correlation or outcome. \\
Function reframing & Whether a reply challenges the view by arguing that the institution or behavior serves a different purpose than the poster assumes. \\
Use-case focus & Whether a reply argues from the specific practical situations in which the disputed thing works especially well or badly. \\
Institutional purpose & How much a reply justifies or criticizes something by appealing to the purpose of the larger system it belongs to. \\
Common-ground signaling & Whether a reply explicitly agrees with part of the poster's position before presenting disagreement. \\
Empathetic tone & Whether a reply validates the poster's feelings or perspective in a supportive way while disagreeing. \\
Motive attack & Whether a reply questions the poster's motives, sincerity, or openness instead of mainly addressing the substance of the claim. \\
Confrontational tone & How much a reply uses insult, scorn, or aggressive language toward the poster or target. \\
Humor deployment & Whether a reply uses jokes, playful phrasing, or irony as part of its persuasive approach. \\
Declared stance framing & Whether a reply explicitly frames itself as devil's advocate, partial agreement, or some other interpretive stance before arguing. \\
Epistemic caution & How much a reply qualifies its claims, notes uncertainty, or avoids overstating what it can prove. \\
Conditional reasoning & Whether a reply limits its claim to specific conditions or exceptions instead of treating the issue as all-or-nothing. \\
Abstraction level & Whether a reply operates mainly at a high-level philosophical or theoretical plane rather than a concrete one. \\
Explanatory correction & Whether a reply mainly persuades by supplying overlooked background information or correcting a misunderstanding. \\
Personal remedy suggestion & Whether a reply addresses the issue by recommending a practical step the poster could take in their own life. \\
Audience reframing & Whether a reply changes the evaluation of the thing in question by showing it was made for a different audience or goal than the poster assumes. \\
Scope challenge & Whether a reply argues that the poster's principle is too narrow to handle the full range of relevant cases. \\
Temporal horizon & Whether a reply evaluates the claim by emphasizing long-term effects rather than immediate ones, or vice versa. \\
Norm distinction & Whether a reply distinguishes between what is logically valid, socially useful, morally justified, or legally allowed rather than treating them as the same question. \\
Tradition appeal & Whether a reply supports or critiques a position by invoking established custom, continuity, or inherited social practice. \\
Privacy sensitivity & Whether a reply foregrounds the need to protect personal vulnerability, anonymity, or intimate information when evaluating the issue. \\
Collective agency framing & Whether a reply explains outcomes as the aggregate result of many individuals' choices rather than the intent of a single actor. \\
Resource referral & Whether a reply mainly tries to persuade by handing off a longer external resource instead of making the argument in the reply itself. \\
Form substance split & Whether a reply distinguishes between objection to labels, rhetoric, or presentation and objection to the underlying idea itself. \\
Meta-argumentation & Whether a reply persuades by analyzing the quality or structure of the argument itself rather than primarily debating the topic. \\
\bottomrule
\end{tabularx}
